\documentclass[../../../include/open-logic-section]{subfiles}

\begin{document}

\olfileid[hi]{sth}{spine}{zf}
\olsection{$\Z$ और $\ZF$: एक महत्त्वपूर्ण पड़ाव}

आधारिता के साथ हम एक और महत्त्वपूर्ण पड़ाव पर पहुँचते हैं। हमने सिद्धांतों
$\Zminus$ और $\ZFminus$ पर विचार किया, जिन्हें हमने कुछ सिद्धांतों में से
कुछ ``घटाकर'' प्राप्त कहा था। वह घटाई गई चीज़ आधारिता है। अतः:

\begin{defn}
सिद्धांत $\Z$, $\Zminus$ में आधारिता जोड़ता है। अतः उसके अभिगृहीत विस्तारिता,
सम्मिलन, युग्म, घात समुच्चय, अनंतता, आधारिता और पृथक्करण योजना के सभी
दृष्टांत हैं।

सिद्धांत $\ZF$, $\ZFminus$ में आधारिता जोड़ता है। दूसरे शब्दों में, $\ZF$,
~$\Z$ में प्रतिस्थापन के सभी दृष्टांत जोड़ता है।
\end{defn}

फिर भी आपके मन में एक प्रश्न आया होगा। यदि $\ZFminus$ में नियमितता आधारिता
के समतुल्य है और नियमितता का औचित्य स्पष्ट है, तो इतना चक्कर लगाकर नियमितता
के बजाय आधारिता को अपना मूल अभिगृहीत क्यों मानें?

ऐतिहासिक कारणों (किसने क्या और कब सूत्रित किया) को अलग रखें, तो मूल कारण यह
है कि आधारिता को~$V_\alpha$ की परिभाषा का उपयोग किए बिना प्रस्तुत किया जा
सकता है। वह परिभाषा \olref[recursion]{sec} के पूरे काम पर निर्भर थी: उसका
औचित्य दिखाने के लिए हमें परिमितेतर पुनरावर्तन सिद्ध करना पड़ा। किंतु
परिमितेतर पुनरावर्तन के हमारे प्रमाण ने \emph{प्रतिस्थापन} का उपयोग किया।
अतः जहाँ आधारिता और नियमितता $\ZFminus$ के सापेक्ष समतुल्य हैं, वे
$\Zminus$ के सापेक्ष समतुल्य नहीं हैं।

वास्तव में बात इस सरल टिप्पणी से कहीं अधिक गंभीर है। यद्यपि यह इस पुस्तक के
दायरे से बहुत बाहर जाता है, पता चलता है कि $\Zminus$ और $\Z$ दोनों
$V_\alpha$ को परिभाषित करने के लिए अत्यंत दुर्बल हैं। अतः यदि आप केवल $\Z$
में काम कर रहे हैं, तो नियमितता (हमारे सूत्रित रूप में) का \emph{अर्थ} तक
नहीं बनता। इसी कारण हमारा आधिकारिक अभिगृहीत नियमितता के बजाय आधारिता है।

अब से हम बिना और टिप्पणी के $\ZF$ में काम करेंगे (जब तक अन्यथा न कहा जाए)।

\end{document}
